\begin{abstract}
\vspace{-3mm}
Recent advances in large language models (LLMs) have established conversational text-to-SQL as a practical interface between users and databases, often involving multiple turns of clarification and revision. However, existing benchmarks primarily evaluate execution accuracy, leaving the unfolding and shifting of user intent across turns largely uncovered. To address this, we introduce \textbf{\textsc{TIDE-Bench}}, a benchmark for conversa\textbf{\textsc{ti}}onal text-to-SQL under chain ambiguity and intent \textbf{\textsc{d}}rift \textbf{\textsc{e}}valuation, targeting two recurring patterns: \textbf{chain ambiguity}, where an underspecified question triggers layered clarification with conditional dependencies, and \textbf{intent drift}, where the user retracts and replaces a previously committed request element. Built on 514 anchor SQLs from BIRD, \textsc{TIDE-Bench} comprises 1{,}542 samples and introduces dedicated metrics for chain identification and drift recognition-resolution beyond execution accuracy. Evaluating 12 advanced LLMs reveals a persistent chain identification bottleneck unaffected by clarification frequency, a wide drift recognition-resolution gap, and overlap between failure modes when jointly activated.
The corresponding code of \textsc{TIDE-Bench} is released for further research\footnote{\href{https://github.com/Rcrossmeister/TIDE-Bench}{https://github.com/Rcrossmeister/TIDE-Bench}}. 
\end{abstract}
